\section{Глава~\ref{sec:dishbrain}. DishBrain}

Устройство стенда и результаты игры взяты из одной экспериментальной работы и её продолжения; формулы шума и дрейфа к хорошим настройкам выведены в главе и проверены моделью книги, а механизм обучения в чашке не установлен.

{\footnotesize
\setlength{\tabcolsep}{3pt}\renewcommand{\arraystretch}{1.15}
\begin{longtable}{>{\raggedright}p{0.5cm}>{\raggedright}p{4.0cm}>{\raggedright}p{5.6cm}>{\raggedright}p{2.3cm}>{\raggedright\arraybackslash}p{1.7cm}}
\toprule
№ & Утверждение & Опыт и что измерено & Источник & Статус \\
\midrule\endfirsthead
\toprule
№ & Утверждение & Опыт и что измерено & Источник & Статус \\
\midrule\endhead
1 & Стенд: около $800\,000$ клеток коры мыши или около $10^6$ клеток человека на чипе; $26\,000$ электродов на $8$~мм$^2$, запись до $1024$, стимуляция через $8$ & Методы работы: чип MaxOne, $26\,000$ платиновых электродов на $8$~мм$^2$, запись $1024$ каналов по $20\,000$ отсчётов в секунду; стимулировать технически можно до $32$ электродов, независимо удалось $8$; мышиные культуры использовали примерно с $10$~суток & \cite{Kagan2022} & измерено \\
2 & Петля с тактом $10$~мс: импульсы двигательных областей двигают ракетку (рис.~\ref{fig:dishloop}) & Импульсы считаются окнами по $10$~мс; задержка от импульса до ответного стимула около $5$~мс, её задаёт буфер аппаратуры & \cite{Kagan2022} & измерено \\
3 & Вход: код местом (какой из $8$ электродов) и частотой $4$--$40$~имп/с & В основном опыте частота стимуляции растёт линейно от $4$~Гц, когда мяч у дальней стенки, до $40$~Гц у ракетки; амплитуда импульсов мяча --- $75$~мВ; импульсы двухфазные прямоугольные & \cite{Kagan2022} & измерено \\
4 & Выход $\Delta p=c\,(n_\uparrow-n_\downarrow)$~\eqref{eq:dishreadout}, шум счёта за окно $1/\sqrt{RT}$ & Правило разности двух областей описано в работе (с весами областей по активности), точная формула не дана. Оценка шума --- следствие пуассоновского счёта, выведена в главе & \cite{Kagan2022}; вывод в главе & теория \\
5 & Учитель: после удара --- предсказуемый стимул $100$~имп/с, $0.1$~с; после промаха --- непредсказуемый, $150$~мВ, $5$~имп/с, $4$~с & Методы: после удара $75$~мВ, $100$~Гц, $100$~мс на все $8$ электродов сразу; после промаха $150$~мВ, $5$~Гц в случайные электроды в случайные моменты в течение $4$~с. Дофамина в культуре нет & \cite{Kagan2022} & измерено \\
6 & Сеть действует так, чтобы её вход стал предсказуемым & Принцип свободной энергии --- теоретическое предложение, из которого авторы выводили протокол; опыт с ним согласуется, но не отличает его от более простых механизмов (строки~7--8) & \cite{Friston2010,Kagan2022} & гипотеза \\
7 & Если неудача встряхивает сеть, а удача нет, время в настройке $\rho\propto1/P_{\text{пр}}$~\eqref{eq:dishdensity} (утверждение~\ref{prop:dishscramble}) & Следует из баланса потоков в марковской цепи с симметричными толчками, выведено в главе. Прямой проверки на культуре нет & вывод в главе & теория \\
8 & Модель «перемешивание при промахе» учится: доля ударов $0.21\to0.38$; при толчке после каждого розыгрыша $\approx0.12$, без толчков $0.19$ (рис.~\ref{fig:dishmodel}) & Расчёт, скрипт \texttt{models/dishbrain\_model.py}, $40$ моделей-культур по $2000$ розыгрышей: $0.21\to0.38$; $0.13\to0.11$; $0.19\to0.19$ & --- & модель \\
9 & Серии отбитых мячей удлиняются только у живых культур с обратной связью; контроли не учатся & $399$ сеансов по $20$~мин: $9$ культур мыши ($101$ сеанс), $11$ человека ($138$), $6$ чипов со средой ($80$), $20$ культур в покое ($42$), $3$ симуляции ($38$). Средняя длина розыгрыша за последние $15$~мин больше, чем за первые $5$, только у живых культур; у клеток, не являющихся нейронами (HEK293T), и у чипов со средой эффекта нет & \cite{Kagan2022} & измерено \\
10 & Без обратной связи и при тишине вместо стимула улучшения нет & $486$ сеансов за $3$ дня: «стимул» ($n=27$), «тишина» ($17$), «без обратной связи» ($15$). Рост во времени значим только в условии «стимул». Оговорка: в конце сеанса условие «тишина» всё же превосходило «без обратной связи», хоть и с меньшим эффектом & \cite{Kagan2022} & измерено \\
11 & Эффект невелик; учится ли сеть надолго --- открытый вопрос & Внутри сеанса улучшение надёжно, но, по словам самих авторов, обучение между сеансами за несколько дней устойчиво не наблюдалось & \cite{Kagan2022} & измерено \\
12 & Во время игры сеть приближается к критическому режиму, и чем ближе, тем лучше игра & Анализ лавин активности тех же культур: структурированный вход приближает сеть к критическому режиму, качество игры коррелирует с близостью к нему; но одной критичности без сведений о последствиях действий для обучения мало & \cite{Habibollahi2023} & измерено \\
13 & «Разумность» культуры не показана & Ответная статья тридцати исследователей: изменение поведения в замкнутой петле не доказывает ни разума, ни ощущений; опыта, который это показывал бы, нет & \cite{Balci2023} & гипотеза \\
14 & Предложенный четвёртый контроль: предсказуемый стимул после промаха той же длительности и энергии (раздел~\ref{sec:dishspec}) & Такой опыт не поставлен; это предложение книги & --- & гипотеза \\
\bottomrule
\end{longtable}}
